\section{Introduction}

Transformer-based Large Language Models (LLMs) are exceptionally efficient containers of information, but remain infamously opaque to classical modular decomposition. Although modern LLMs are often described as over-parameterized, the number of bits they use to represent their parameters is small compared to the information content of their training datasets; even more so under quantization/reduced-precision frameworks. Standard language-model training minimizes prediction loss over the contents of a dataset. In this sense, learning can be viewed as a form of lossy information transfer, shaping LLMs into the most efficient compressed representations of their training corpora with respect to their training objective.
\newline
This heavily bit-constrained regime also makes tracing the writing of new information into a model, via training, fundamentally nontrivial. Even parameter-efficient post-training methods such as Low-Rank Adaptation (LoRA) \citep{hu2022lora}, despite being modular in parameter space, still induce behavioral changes through distributed interactions with the frozen model’s internal representations. When finetuning pretrained models, parameter-efficient adapters provide a clean setting for studying this process, as they isolate the written information to a parameterized update whose effective bit budget can be directly controlled.
\newline
In this work, we adopt an information-theoretic view of finetuning. We treat the task-relevant information supplied by the finetuning data as a source, and the finetuning process as an encoder of this information. The adapter in which this information is written can be thought of as a finite-rate message; as it carries a lossy parameter-space encoding of the finetuned behavior. During inference, the composition of the frozen base model and the adapter produces the acquired behavior, i.e., the reconstruction of the information supplied during training.
% --- Palette: figures4papers (paper/plots/style.py), same meanings as the plots:
% blue = the adapter / message, gray = frozen base, red = damage/distortion.
\definecolor{srcFill}{HTML}{EEF7EE}
\definecolor{srcDraw}{HTML}{8BCF8B}
\definecolor{srcIcon}{HTML}{3F8F3F}
\definecolor{encFill}{HTML}{F5EAF3}
\definecolor{encDraw}{HTML}{C28FB9}
\definecolor{encIcon}{HTML}{9A4D8E}
\definecolor{msgFill}{HTML}{E6EEF7}
\definecolor{msgDraw}{HTML}{3775BA}
\definecolor{msgIcon}{HTML}{0F4D92}
\definecolor{baseFill}{HTML}{F2F2F2}
\definecolor{decDraw}{HTML}{4D4D4D}
\definecolor{baseIcon}{HTML}{767676}
\definecolor{adpFill}{HTML}{E6EEF7}
\definecolor{adpIcon}{HTML}{0F4D92}
\definecolor{outFill}{HTML}{E5F2F3}
\definecolor{outDraw}{HTML}{42949E}
\definecolor{outIcon}{HTML}{2C6F77}
\definecolor{distDraw}{HTML}{B64342}

\begin{figure}[t]
\centering
\resizebox{\linewidth}{!}{%
\begin{tikzpicture}[
  font=\normalsize,
  >={Latex[length=2.7mm,width=2.0mm]},
  flow/.style={->,line width=1.15pt,draw=decDraw},
  panel/.style={
    rounded corners=2.5pt,
    line width=1.0pt,
    minimum width=3.15cm,
    minimum height=2.85cm,
    inner sep=0pt
  },
  role/.style={font=\bfseries\large},
  main/.style={font=\large,align=center,text=black!92},
  sub/.style={font=\itshape\normalsize}
]

% --- Main panels ---------------------------------------------------
\node[panel,draw=srcDraw,fill=srcFill] (data) at (0,0) {};
\node[panel,draw=encDraw,fill=encFill] (training) at (4.05,0) {};
\node[panel,draw=msgDraw,fill=msgFill] (adapter) at (8.10,0) {};
\node[minimum width=4.80cm,minimum height=2.85cm,inner sep=0pt] (decoder) at (13.10,0) {};
\node[panel,draw=outDraw,fill=outFill] (behavior) at (18.15,0) {};

% --- SOURCE icon: clean stacked documents --------------------------
\begin{scope}[
    shift={($(data.center)+(0,0.64)$)},
    draw=srcIcon,
    line width=1.15pt,
    scale=1.18
]

% back sheet
\draw[
    fill=srcFill,
    rounded corners=1.2pt
]
(-0.02,-0.28) rectangle (0.34,0.30);

% middle sheet
\draw[
    fill=srcFill,
    rounded corners=1.2pt
]
(-0.12,-0.20) rectangle (0.24,0.38);

% front sheet
\draw[
    fill=srcFill,
    rounded corners=1.2pt
]
(-0.22,-0.12) rectangle (0.14,0.46);

% text lines on front sheet only
\draw[line width=.9pt] (-0.14,0.26) -- (0.06,0.26);
\draw[line width=.9pt] (-0.14,0.13) -- (0.06,0.13);
\draw[line width=.9pt] (-0.14,0.00) -- (0.02,0.00);

\end{scope}
\node[main] at ($(data.center)+(0,-0.70)$) {Fine-tuning\\data};

% --- ENCODER icon --------------------------------------------------
\begin{scope}[shift={($(training.center)+(0,0.66)$)},scale=1.18]
  \foreach \a in {0,45,...,315}{\begin{scope}[rotate=\a]\fill[encIcon] (-0.08,0.29) rectangle (0.08,0.43);\end{scope}}
  \fill[encIcon] (0,0) circle (0.34);
  \fill[encFill] (0,0) circle (0.13);
\end{scope}
\node[main] at ($(training.center)+(0,-0.62)$) {Fine-tuning};

% --- MESSAGE adapter glyph ----------------------------------------
\begin{scope}[shift={($(adapter.center)+(0,0.66)$)},draw=msgIcon,line width=1.25pt,scale=1.18]
  \draw[fill=msgFill,rounded corners=2.6pt] (-0.34,-0.31) rectangle (0.40,0.31);
  \foreach \yy in {-0.18,0,0.18}{\draw (-0.50,\yy)--(-0.34,\yy);}
  \draw (0.02,-0.17)--(0.02,0.17);
  \draw (0.14,-0.17)--(0.14,0.17);
  \draw (0.26,-0.17)--(0.26,0.17);
\end{scope}
\node[main] at ($(adapter.center)+(0,-0.70)$) {Learned\\adapter};

% --- DECODER -------------------------------------------------------
\begin{scope}
  \clip[rounded corners=2.5pt] (decoder.south west) rectangle (decoder.north east);
  \fill[baseFill] (decoder.south west) rectangle (decoder.center |- decoder.north);
  \fill[adpFill] (decoder.center |- decoder.south) rectangle (decoder.north east);
\end{scope}
\draw[draw=decDraw,line width=1.0pt,rounded corners=2.5pt]
  (decoder.south west) rectangle (decoder.north east);

% docking divider
\draw[dashed,dash pattern=on 3pt off 2.4pt,draw=black!34,line width=.72pt]
  ($(decoder.center)+(0,1.425)$)--($(decoder.center)+(0,0.38)$);
\draw[draw=black!34,line width=.72pt]
  ($(decoder.center)+(0,0.38)$) arc[start angle=90,end angle=-90,radius=.20cm];
\draw[dashed,dash pattern=on 3pt off 2.4pt,draw=black!34,line width=.72pt]
  ($(decoder.center)+(0,-0.38)$)--($(decoder.center)+(0,-1.425)$);

% base-model network icon
% --- Base-model icon: compact layered neural network ----------------
\begin{scope}[
    shift={($(decoder.west)!0.5!(decoder.center)+(0,0.49)$)},
    x=0.34cm,
    y=0.27cm
]

% connections: input -> hidden
\foreach \yi in {-1,0,1}{
    \foreach \yh in {-1.35,-0.45,0.45,1.35}{
        \draw[
            draw=baseIcon!55,
            line width=0.45pt
        ] (-1,\yi) -- (0,\yh);
    }
}

% connections: hidden -> output
\foreach \yh in {-1.35,-0.45,0.45,1.35}{
    \foreach \yo in {-1,0,1}{
        \draw[
            draw=baseIcon!55,
            line width=0.45pt
        ] (0,\yh) -- (1,\yo);
    }
}

% nodes
\foreach \y in {-1,0,1}{
    \filldraw[
        fill=baseFill,
        draw=baseIcon,
        line width=0.75pt
    ] (-1,\y) circle (0.105cm);

    \filldraw[
        fill=baseFill,
        draw=baseIcon,
        line width=0.75pt
    ] (1,\y) circle (0.105cm);
}

\foreach \y in {-1.35,-0.45,0.45,1.35}{
    \filldraw[
        fill=baseFill,
        draw=baseIcon,
        line width=0.75pt
    ] (0,\y) circle (0.105cm);
}

\end{scope}
\node[font=\large,text=black!92] at ($(decoder.west)!0.5!(decoder.center)+(0,-0.48)$) {Base model};
\node[sub,text=baseIcon] at ($(decoder.west)!0.5!(decoder.center)+(0,-0.91)$) {frozen};

% matching trained-adapter glyph
\begin{scope}[shift={($(decoder.center)!0.5!(decoder.east)+(0,0.70)$)},draw=adpIcon,line width=1.25pt,scale=1.18]
  \draw[fill=adpFill,rounded corners=2.6pt] (-0.34,-0.31) rectangle (0.40,0.31);
  \foreach \yy in {-0.18,0,0.18}{\draw (-0.50,\yy)--(-0.34,\yy);}
  \draw (0.02,-0.17)--(0.02,0.17);
  \draw (0.14,-0.17)--(0.14,0.17);
  \draw (0.26,-0.17)--(0.26,0.17);
\end{scope}
\node[font=\large,text=black!92] at ($(decoder.center)!0.5!(decoder.east)+(0,-0.48)$) {Adapter};
\node[sub,text=adpIcon] at ($(decoder.center)!0.5!(decoder.east)+(0,-0.91)$) {trained};

% --- RECONSTRUCTION icon ------------------------------------------
\begin{scope}[shift={($(behavior.center)+(0,0.66)$)},draw=outIcon,line width=1.25pt,scale=1.18]
  \draw (0,0) circle (.34);
  \draw (0,0) circle (.21);
  \draw (0,0) circle (.07);
  \draw[->] (.39,.39)--(.0,.0);
\end{scope}
\node[main] at ($(behavior.center)+(0,-0.70)$) {Adapted\\behavior};

% --- Stage headings ------------------------------------------------
\node[role,text=srcIcon] at ($(data.north)+(0,1.00)$) {SOURCE};
\node[role,text=encIcon] at ($(training.north)+(0,1.00)$) {ENCODER};
\node[role,text=msgIcon] at ($(adapter.north)+(0,1.00)$) {MESSAGE};
\node[role,text=decDraw] at ($(decoder.north)+(0,1.00)$) {DECODER};
\node[role,text=outIcon] at ($(behavior.north)+(0,1.00)$) {RECONSTRUCTION};

% heading brackets
\foreach \node/\col in {data/srcIcon,training/encIcon,adapter/msgIcon,decoder/decDraw,behavior/outIcon}{
  \draw[\col,line width=.75pt,rounded corners=2pt]
    ($ (\node.north west)+(0,0.38) $)
    -- ($ (\node.north west)+(0,0.50) $)
    -- ($ (\node.north east)+(0,0.50) $)
    -- ($ (\node.north east)+(0,0.38) $);
}

% --- Main flow -----------------------------------------------------
\draw[flow] (data.east)--(training.west);
\draw[flow] (training.east)--(adapter.west);
\draw[flow] (adapter.east)--(decoder.west);
\draw[flow] (decoder.east)--(behavior.west);

% --- Rate brace ----------------------------------------------------
\draw[decorate,decoration={brace,mirror,amplitude=4.4pt},draw=msgIcon,line width=.8pt]
  ($(adapter.south west)+(0,-.12)$)--($(adapter.south east)+(0,-.12)$);
\node[font=\large,text=msgIcon] (rate) at ($(adapter.south)+(0,-.62)$) {rate $R$};

% --- Behavioral distortion ----------------------------------------
\coordinate (dbot) at ($(adapter.south)+(0,-1.46)$);
\draw[->,dash pattern=on 4pt off 3pt,line width=.9pt,draw=distDraw]
  (data.south)
  -- (data.south |- dbot)
  -- node[pos=.53,below=5pt,font=\large,text=distDraw,fill=white,inner xsep=2pt]
     {behavioral distortion $D$}
  (behavior.south |- dbot)
  -- (behavior.south);

\end{tikzpicture}%
}

\caption{Information-theoretic view of parameter-efficient fine-tuning. $R$ is the number of bits used by the adapter, $D$ measures failure to preserve the behavior induced by fine-tuning.}
\label{fig:information_view}
\end{figure}

Under this analogy, illustrated in Figure \ \ref{fig:information_view}, we can also reuse the concepts of rate $R$, as the number of bits used to represent the adapter; while distortion $D$ is a measure of the degradation of the original behavior upon reconstruction. Put together, we establish a rate-distortion relation as the trade-off between the bit budget of the adapter and the preservation of the behavior induced by finetuning. Finally, for a fixed distortion tolerance $D^\star$, we can define the behavioral rate $R^\star$ as the smallest bit budget allowing for the reconstruction of an adapter whose behavior satisfies $D^\star$. In practice, we instantiate $R$ by training an adapter once, compressing it and measuring the resulting file size, then evaluate its performance after decoding back into the frozen model. This way, we can measure $D$ behaviorally, rather than using parameter-space proxies such as the adapter's weight reconstruction error. $R^\star(D^\star)$ is thus a measure of the smallest tested file size, and can be measured in compressed adapter file bits per initial parameter.
\newline
We first show that ordinary proxies do not predict $R^\star$: neither rank, parameter count, weight reconstruction error, update magnitude nor the size of the behavioral gain sets it, and rank and precision are not interchangeable ways of spending bits. The rate is receiver-relative: the same datasets need different rates on different frozen models, sometimes in reversed order. Correction spectra measured on the data explain these differences within a receiver but not across receivers; the trained adapter's own spectrum does better.
\newline
Most importantly, we find that the information written by fine-tuning has an uneven impact. In structured tasks like reasoning-based problem-solving, training an adapter deliberately on corrupted datasets can severely damage task performance over a base model. However, by compressing the adapter, we are consistently able to recover most of the lost behavior, and even surpass the frozen base model across several types of tasks, model sizes and family. The damage sits in the adapter's leading singular directions: coarse coding works mostly by shrinking them, and removing them at full precision works as well. The compressed adapter also keeps the frozen model's accuracy on unrelated tasks, where a clean fine-tune loses it.
\newline
Together, these results suggest that finetuning information is structured into components of sharply different behavioral value, and that this value depends on both the receiver and the adapter carrier.

% svd results on corrrupted work.